% Generated by scripts/assemble.py; corrections belong in source/corrections.json.
\examhead{140}{华中师范大学 2026 年数学分析试卷(回忆版)}{133}
% p134-b19; source PDF page 134
\blockanchor{p134-b19}
(注. 题目来自考生回忆, 部分表述或数据可能有误)\par
% p134-b20; source PDF page 134
\par\Needspace{5\baselineskip}
\blockanchor{p134-b20}
\majorquestion{一}计算题.\par
% p134-b21; source PDF page 134
\blockanchor{p134-b21}
(1) 求极限  $\lim_{x\to0}x^{-4}\left(\cos x-e^{-\frac{x^{2}}{2}}\right)$\par
% p134-b22; source PDF page 134
\blockanchor{p134-b22}
(2) 求由 $z = x^{2} + y^{2}$, $z = x + y$ 所围成的立体图形的体积.\par
% p134-b23; source PDF page 134
\par\Needspace{6\baselineskip}
\blockanchor{p134-b23}
(3) 计算曲面积分\par
\nopagebreak[4]
% p134-b24; source PDF page 134
\blockanchor{p134-b24}
\displayformula{\iint_ {S} x z ^ {2} \mathrm{d} y \mathrm{d} z + (x ^ {2} y - z ^ {3}) \mathrm{d} z \mathrm{d} x + (2 x y + y ^ {2} z) \mathrm{d} x \mathrm{d} y.}
% p134-b25; source PDF page 134
\blockanchor{p134-b25}
其中 $S$ 是由 $z = \sqrt{a^2 - x^2 - y^2} (a > 0)$ 与 $z = 0$ 所围立体的表面, 取外侧为正向.\par
% p134-b26; source PDF page 134
\par\Needspace{5\baselineskip}
\blockanchor{p134-b26}
\majorquestion{二}解答如下问题:\par
% p135-b00; source PDF page 135
\blockanchor{p135-b00}
(1) 设 $L$ 为常数, $\{y_{n}\}$ 是严格递增的正无穷大数列, 且与数列 $\{x_{n}\}$ 一起满足 $\lim_{n\to \infty}\frac{x_{n + 1} - x_n}{y_{n + 1} - y_n} = L$. 证明: $\lim_{n\to \infty}\frac{x_n}{y_n} = L$.\par
% p135-b01; source PDF page 135
\blockanchor{p135-b01}
(2) 设  $a_{1} = \sin a > 0, a_{n+1} = \sin a_{n} (n = 1, 2, \cdots)$ ，证明： $\lim_{n \to \infty} \frac{1}{na_{n}^{2}} = \frac{1}{3}$ .\par
% p135-b02; source PDF page 135
\par\Needspace{5\baselineskip}
\blockanchor{p135-b02}
\majorquestion{三}设 $f(x)$ 在 $(a, b)$ 上可导，$\lim_{x \to a^{+}} f'(x)$ 与 $\lim_{x \to b^{-}} f'(x)$ 均存在。证明：\par
% p135-b03; source PDF page 135
\blockanchor{p135-b03}
(1) $f(a + 0)$ 和 $f(b - 0)$ 都存在.\par
% p135-b04; source PDF page 135
\blockanchor{p135-b04}
(2) $f(x)$ 在 $(a, b)$ 上有界，且一致连续.\par
% p135-b05; source PDF page 135
\par\Needspace{5\baselineskip}
\par\Needspace{6\baselineskip}
\blockanchor{p135-b05}
\majorquestion{四}设 $f(x)$ 在 $I = (0, +\infty)$ 上可导, 且对任意的 $x_{1}, x_{2} \in I$ 以及 $\lambda_{1}, \lambda_{2} \in (0, 1)$, 有\par
\nopagebreak[4]
% p135-b06; source PDF page 135
\blockanchor{p135-b06}
\displayformula{f (\lambda_ {1} x _ {1} + \lambda_ {2} x _ {2}) \leq \lambda_ {1} f (x _ {1}) + \lambda_ {2} f (x _ {2}).}
% p135-b07; source PDF page 135
\blockanchor{p135-b07}
其中 $\lambda_1 + \lambda_2 = 1$。假设 $x \in (0, +\infty)$ 时，有 $f(x) = x^2 + x^2\varepsilon(x)$，其中 $\lim_{x \to +\infty} \varepsilon(x) = 0$。证明：\par
% p135-b08; source PDF page 135
\blockanchor{p135-b08}
(1) 对任意的 $x \in (0, +\infty)$, 有 $\frac{f(x) - f(x - h)}{h} \leq f'(x) \leq \frac{f(x + h) - f(x)}{h}, 0 < h < \frac{x}{2}$.\par
% p135-b09; source PDF page 135
\blockanchor{p135-b09}
(2) 对任意的 $\eta > 0$, 当 $x$ 充分大时, 有 $2x - h - \frac{\eta}{h} x^2 \leq f'(x) \leq 2x + h + \frac{\eta}{h} x^2$, $0 < h < \frac{x}{2}$.\par
% p135-b10; source PDF page 135
\par\Needspace{5\baselineskip}
\blockanchor{p135-b10}
\majorquestion{五}设 $f(x)$ 为有界闭区间 $[a, b]$ 上的函数, 且对任意的 $x_0 \in [a, b]$, 有 $\lim_{x \to x_0} f(x) = 0$. 证明:\par
% p135-b11; source PDF page 135
\blockanchor{p135-b11}
(1) $f(x)$ 在 $[a, b]$ 上有界.\par
% p135-b12; source PDF page 135
\blockanchor{p135-b12}
(2) 对任意的 $\varepsilon > 0$, 在 $[a, b]$ 上至多存在有限个点 $x$, 使得 $|f(x)| \geq \varepsilon$.\par
% p135-b13; source PDF page 135
\blockanchor{p135-b13}
(3) $f(x)$ 在 $[a, b]$ 上可积, 且 $\int_{a}^{b} f(x) \mathrm{d}x = 0$.\par
% p135-b14; source PDF page 135
\par\Needspace{5\baselineskip}
\blockanchor{p135-b14}
\majorquestion{六}判断 $\sum_{n=3}^{\infty} \frac{1}{n^{p} (\ln n)^{q} (\ln \ln n)^{r}}$ 的敛散性.\par
% p135-b15; source PDF page 135
\par\Needspace{5\baselineskip}
\blockanchor{p135-b15}
\majorquestion{七}设函数列 $\{f_n(x)\}$ 与 $\{g_n(x)\}$ 在区间 $I$ 上分别一致收敛于 $f(x)$ 与 $g(x)$, 且存在正数列 $\{M_n\}$, 使得当 $x \in I, n = 1,2,\dots$ 时, 有 $|f_n(x)| \leq M_n$, $|g_n(x)| \leq M_n$. 证明: $\{f_n(x)g_n(x)\}$ 在 $I$ 上一致收敛.\par
% p135-b16; source PDF page 135
\par\Needspace{5\baselineskip}
\blockanchor{p135-b16}
\majorquestion{八}已知 $f(x)$ 是以 $2\pi$ 为周期的可积函数, 通过 $f(x)$ 的傅里叶系数表示 $f(x + h)$ 的傅里叶系数.\par
% p135-b17; source PDF page 135
\par\Needspace{5\baselineskip}
\blockanchor{p135-b17}
\majorquestion{九}设 $f(x, y) = |x - y| g(x, y)$, 其中 $g(x, y)$ 在 (0, 0) 的某邻域内连续. 问:\par
% p135-b18; source PDF page 135
\blockanchor{p135-b18}
(1) 在什么条件下, $f_{x}(0,0), f_{y}(0,0)$ 存在?\par
% p135-b19; source PDF page 135
\blockanchor{p135-b19}
(2) 在什么条件下, $f(x,y)$ 在 $(0,0)$ 处可微?\par
% p135-b20; source PDF page 135
\par\Needspace{5\baselineskip}
\blockanchor{p135-b20}
\majorquestion{十}设函数 $f(x)$ 在 $(0, +\infty)$ 上有界且连续, 若对任意的 $c \in \mathbb{R}$, 函数 $f(x) - c$ 在 $(0, +\infty)$ 上至多有有限个零点, 证明: $\lim_{x \to +\infty} f(x)$ 存在.\par
